\documentclass[12pt,a4paper]{article}
\usepackage[margin=25mm, headheight=15pt, headsep=10pt]{geometry}
\usepackage{fontspec}
\usepackage[table]{xcolor} % Note: table option is required for rowcolors
\usepackage{titlesec}
\usepackage{tocloft}
\usepackage{fancyhdr}
\usepackage{booktabs}
\usepackage{tcolorbox}
\tcbuselibrary{skins, breakable}
\usepackage[colorlinks=true, linkcolor=emerald, urlcolor=emerald, bookmarks=true]{hyperref}

\definecolor{emerald}{RGB}{0,102,51}
\definecolor{darkred}{RGB}{139,0,0}
\definecolor{lightgray}{RGB}{245,245,245}

\usepackage[nil]{babel}
\babelprovide[import, main]{bengali}
\babelprovide[import]{arabic}
\babelfont{rm}[Renderer=HarfBuzz,Scale=1.0]{Noto Serif Bengali}
\babelfont{sf}[Renderer=HarfBuzz]{Noto Sans Bengali}
\babelfont{tt}[Renderer=HarfBuzz]{Noto Sans Bengali}
\babelfont[arabic]{rm}[Renderer=HarfBuzz,Scale=1.1]{Noto Naskh Arabic}

% Section styling
\titleformat{\section}{\Large\bfseries\color{emerald}}{\thesection.}{0.5em}{}
\titleformat{\subsection}{\large\bfseries\color{emerald!85!black}}{\thesubsection.}{0.5em}{}
\titleformat{\subsubsection}{\normalsize\bfseries\color{emerald!70!black}}{\thesubsubsection.}{0.5em}{}
\titlespacing*{\section}{0pt}{18pt}{8pt}

% Fancy Header/Footer setup
\pagestyle{fancy}
\fancyhf{} % clear all header and footer fields
\fancyhead[L]{\color{emerald}\small\textbf{\nouppercase{\leftmark}}}
\fancyfoot[L]{\scriptsize\color{black!50}{\lat{AI}}-সংকলিত, আলেম-যাচাইকৃত নয়}
\fancyfoot[C]{\color{emerald!70!black}\thepage}
\renewcommand{\headrulewidth}{0.4pt}
\renewcommand{\headrule}{\hbox to\headwidth{\color{emerald}\leaders\hrule height \headrulewidth\hfill}}
\renewcommand{\footrulewidth}{0pt}

% Custom boxes
% 1. Ayat/Hadith Box
\newtcolorbox{ayatbox}[1][]{
    enhanced, breakable, 
    colback=emerald!5, 
    colframe=emerald, 
    boxrule=0.5pt, 
    arc=3pt, 
    left=10pt, right=10pt, top=10pt, bottom=10pt,
    #1
}

% 2. Quote/Translation Box
\newtcolorbox{quotebox}[1][]{
    enhanced, breakable,
    frame hidden,
    colback=lightgray,
    borderline west={3pt}{0pt}{emerald},
    left=15pt, right=10pt, top=10pt, bottom=10pt,
    sharp corners,
    #1
}

% 3. AI-generated notice box
\newtcolorbox{warnbox}[1][]{
    enhanced, breakable,
    colback=red!4,
    colframe=darkred,
    boxrule=0.8pt,
    arc=3pt,
    left=10pt, right=10pt, top=10pt, bottom=10pt,
    #1
}

% Commands
\newcommand{\arab}[1]{\foreignlanguage{arabic}{#1}}
\newcommand{\kib}[1]{\textcolor{emerald}{\textbf{#1}}}

% Latin (English) fallback: Noto Serif Bengali has NO Latin glyphs
\newfontfamily\latfont[Renderer=HarfBuzz]{Noto Serif}
\newcommand{\lat}[1]{{\latfont #1}}

\setlength{\parskip}{5pt}
\setlength{\parindent}{0pt}

% Fix for missing bullet in Noto Serif Bengali
\renewcommand{\labelitemi}{$\bullet$}



\begin{document}
\begin{center}{\Large\bfseries\color{emerald} 07-আবু-দুজানা-যুদ্ধক্ষেত্রে-অহংকার}\end{center}
\vspace{1em}
\section*{ঘটনা ৭: আবু দুজানা (রাঃ) — যুদ্ধক্ষেত্রে অহংকারের ব্যতিক্রম}

\begin{warnbox}
 \textbf{গুরুত্বপূর্ণ নোটিশ (\lat{Important} \lat{Notice}):} এই কন্টেন্টটি \lat{AI} (কৃত্রিম বুদ্ধিমত্তা) দিয়ে প্রস্তুত ও সংকলিত। \lat{AI} কঠোর সূত্র-মান নীতি মেনে শুধু নির্ভরযোগ্য উৎস থেকে তথ্য সংগ্রহ করেছে — কেবল সহীহ/হাসান হাদিস ও সহীহ সনদের আসার রাখা হয়েছে, দুর্বল সনদ বাদ দেওয়া হয়েছে। তবে এটি কোনো আলেম বা মুহাদ্দিস কর্তৃক যাচাইকৃত নয়।
\end{warnbox}


\medskip\hrule\medskip

\subsection*{১. সূত্র কার্ড}

\begin{table}[h]\centering\small
\begin{tabular}{ll}
বিষয় & বিবরণ \\
\hline
\textbf{বর্ণনাকারী} & খালিদ ইবনে সুলায়মান ইবনে সিমাক — আবু দুজানার নাতি, তাঁর পিতা ও দাদা থেকে \\
\textbf{গ্রন্থ} & আল-মু'জামুল কাবীর (তাবরানী), হাদিস ৬৫০৮; সুনানে ইবনে মাজাহ, হাদিস ৪১৬৪ \\
\textbf{অধ্যায়} & মু'জামুল কাবীর — সাহাবাদের নাম অনুযায়ী; ইবনে মাজাহ — কিতাবুজ-যুহদ \\
\textbf{সনদের মান} & \textbf{হাসান} (ইবনে মাজাহের সূত্রে) \\
\textbf{মূল বিষয়} & অহংকার সর্বত্র নিষিদ্ধ — কিন্তু যুদ্ধক্ষেত্রে সাহসিকতা ভিন্ন বিষয় \\
\hline
\end{tabular}
\end{table}

\medskip\hrule\medskip

\subsection*{২. পটভূমি (কে, কখন, কোথায়, কেন)}

\textbf{আবু দুজানা সিমাক ইবনে খারাশাহ (রাঃ)} — আনসারদের একজন প্রসিদ্ধ যোদ্ধা (\lat{warrior})। উহুদ যুদ্ধ (\lat{battle})-এর আগে নবীজি (সাল্লাল্লাহু আলাইহি ওয়া সাল্লাম) তলোয়ার নিয়ে দাঁড়িয়ে জিজ্ঞেস করেছিলেন: \textbf{\lat{“}কে আমার এই তলোয়ার তার হক আদায় করে নেবে?\lat{”}} আবু দুজানা বলেছিলেন: \textbf{\lat{“}এর হক কী, হে আল্লাহর রাসূল?\lat{”}} নবীজি বললেন: \textbf{\lat{“}যতক্ষণ না এটি বেঁকে যায় (শত্রুর ওপর আঘাত করতে করতে), ততক্ষণ লড়বে।\lat{”}} আবু দুজানা তলোয়ার নিলেন — এবং নিজের মাথায় লাল পাগড়ি বেঁধে যুদ্ধক্ষেত্রে (\lat{battlefield}) নামলেন।

লাল পাগড়ি ছিল আবু দুজানার যুদ্ধের চিহ্ন (\lat{symbol}) — তিনি যুদ্ধে এমন ভঙ্গিতে হাঁটতেন যাকে সাধারণত অহংকারের (\lat{arrogance}) চলা বলা হয়। নবীজি তাকে দেখে একটি গুরুত্বপূর্ণ শিক্ষা (\lat{lesson}) দিলেন।

\medskip\hrule\medskip

\subsection*{৩. পূর্ণ ঘটনার বর্ণনা}

\begin{quote}
\textbf{\lat{“}উহুদ দিবসে আবু দুজানা (রাঃ) মাথায় লাল পাগড়ি বেঁধেছিলেন। রাসূলুল্লাহ (সাল্লাল্লাহু আলাইহি ওয়া সাল্লাম) তাকে দুই সারির মাঝখানে অহংকারী ভঙ্গিতে হাঁটতে দেখলেন এবং বললেন:} \\
\textbf{"এ হলো এমন একটি হাঁটার ভঙ্গি, যা আল্লাহ ঘৃণা করেন — তবে এই স্থান (যুদ্ধক্ষেত্র) ব্যতীত।"\lat{”}}
\end{quote}

অর্থাৎ — সাধারণ জীবনে এই অহংকারী চলাফেরা নিষিদ্ধ; কিন্তু শত্রুর সামনে নিজের সাহস ও শক্তি প্রদর্শন করা — যার ফলে শত্রু ভীত হয় — এটি কিবর নয়, বীরত্বের প্রয়োজনীয় প্রকাশ।

\medskip\hrule\medskip

\subsection*{৪. ধাপে ধাপে বিশ্লেষণ}

\textbf{ধাপ ১ — উহুদ যুদ্ধের প্রস্তুতি:} নবীজি তলোয়ার দিয়ে যুদ্ধের প্রস্তুতি (\lat{preparation}) নিচ্ছেন; আবু দুজানা তলোয়ার নিয়ে যুদ্ধে যাওয়ার প্রতিজ্ঞা (\lat{determination}) করেন।

\textbf{ধাপ ২ — লাল পাগড়ি ও অহংকারী ভঙ্গি:} আবু দুজানা লাল পাগড়ি পরে এমন ভঙ্গিতে (\lat{attitude}) হাঁটছেন — যা সাধারণ অবস্থায় অহংকারের লক্ষণ (\lat{sign})।

\textbf{ধাপ ৩ — নবীজির নজর:} নবীজি (সা.) সাহাবাদের চাল-চলন এমনভাবে নজর রাখতেন (\lat{observe}) যে অহংকারের সামান্য ইঙ্গিতও তাঁকে সতর্ক (\lat{alert}) করত।

\textbf{ধাপ ৪ — শিক্ষা:} নবীজি বললেন — এই হাঁটা আল্লাহ ঘৃণা (\lat{hate}) করেন, তবে যুদ্ধক্ষেত্রে বৈধ (\lat{permissible})। কারণ যুদ্ধে সাহস (\lat{courage}) ও পরাক্রম প্রদর্শন করা ঈমানের শক্তি; অহংকার নয়।

\textbf{ধাপ ৫ — চিরন্তন নীতি:} ইসলামের দৃষ্টিতে অহংকার (কিবর) সর্বত্র নিন্দনীয় (\lat{condemned}) — কিন্তু শত্রুর সামনে দৃপ্ততা (যুদ্ধের পরিস্থিতিতে) আলাদা বিধান (\lat{rule})।

\medskip\hrule\medskip

\subsection*{৫. শব্দের ব্যাখ্যা}

\begin{table}[h]\centering\small
\begin{tabular}{ll}
শব্দ & অর্থ \\
\hline
\textbf{মুখতাল (\arab{مختال})} & অহংকারী ভঙ্গিতে হাঁটা — দম্ভভরে চলা \\
\textbf{ইসাবাহ (\arab{عصابة})} & পাগড়ি/মাথায় বাঁধা কাপড় \\
\textbf{মিশ্যাহ (\arab{مشية})} & হাঁটার ভঙ্গি/চলন \\
\textbf{বায়িদ (\arab{يبغض})} & ঘৃণা করা — আল্লাহর অপছন্দ \\
\textbf{মাওযি' (\arab{موضع})} & স্থান — এখানে যুদ্ধক্ষেত্র \\
\hline
\end{tabular}
\end{table}

\textbf{গুরুত্বপূর্ণ পার্থক্য:}
\begin{itemize}
\item \textbf{কিবর/অহংকার:} নিজেকে বড় ভাবা, মানুষকে হেয় করা — সর্বত্র নিষিদ্ধ।
\item \textbf{যুদ্ধের দৃপ্ততা:} শত্রুর সামনে সাহস ও শক্তি প্রকাশ — বৈধ, এমনকি প্রশংসিত।
\end{itemize}
নবীজি এক লাইনে এই সীমারেখাটি স্পষ্ট করে দিলেন।

\medskip\hrule\medskip

\subsection*{৬. সংশ্লিষ্ট হাদিস ও আয়াত}

\textbf{হাদিস (মুসলিম ৯১):} নবীজি বলেছেন: \textbf{\lat{“}যার হৃদয়ে সরিষার দানার পরিমাণ কিবর থাকবে, সে জান্নাতে প্রবেশ করবে না।\lat{”}}

\textbf{হাদিস (বুখারি ৫৭৮৮):} নবীজি বলেছেন: \textbf{\lat{“}এক ব্যক্তি সুন্দর পোশাক পরে অহংকারে হাঁটছিল; আল্লাহ তাকে পৃথিবীতে ধ্বসিয়ে দিলেন — সে কিয়ামত পর্যন্ত ডুবতে থাকবে।\lat{”}} — অহংকারে হাঁটার শাস্তি।

\textbf{আয়াত (সূরা আল-আনফাল ৮:৬৫):}
\begin{quote}
\textbf{\lat{“}হে নবী! মুমিনদেরকে যুদ্ধের জন্য উৎসাহিত করুন।\lat{”}}
\end{quote}

\textbf{আয়াত (সূরা মুহাম্মদ ৪৭:৪):}
\begin{quote}
\textbf{\lat{“}অতএব তোমরা যখন যুদ্ধে (শত্রুর) মুখোমুখি হও, তখন দৃঢ়ভাবে আঘাত করো।\lat{”}}
\end{quote}

\medskip\hrule\medskip

\subsection*{৭. গভীর শিক্ষা}

\textbf{১. অহংকার — সর্বত্র নিষিদ্ধ (\lat{forbidden})।}
সাধারণ জীবনে অহংকারের কোনো রূপ বৈধ নয় (\lat{valid}) — চলাফেরা, পোশাক, কথায়। কিবরের সংজ্ঞা (\lat{definition}) স্পষ্ট: সত্যকে তুচ্ছ করা ও মানুষকে হেয় করা।

\textbf{২. যুদ্ধক্ষেত্রের ব্যতিক্রম (\lat{exception})।}
শত্রুর সামনে সাহস (\lat{courage}) ও শক্তি প্রদর্শন করা কিবর নয় — এটি ঈমানের প্রয়োজনীয় প্রকাশ (\lat{expression}), যা শত্রুকে ভীত করে। নবীজি এ ব্যতিক্রম স্বীকার করেছেন।

\textbf{৩. নবীজির সাহাবাদের প্রতি নজর (\lat{attention})।}
নবীজি সাহাবাদের সামান্য চলাফেরাও লক্ষ্য করতেন — শিক্ষা (\lat{lesson}) দেয়, মুমিনের প্রতিটি আচরণ-ভঙ্গি নিয়ন্ত্রণে (\lat{discipline}) থাকা উচিত।

\textbf{৪. নিয়তের (\lat{intention}) গুরুত্ব।}
একই হাঁটা — কেউ করে অহংকারে, কেউ করে আল্লাহর পথে শক্তি দেখাতে। নিয়তই আমলের প্রাণ (\lat{soul})। আবু দুজানার নিয়ত ছিল যুদ্ধের দৃপ্ততা; তাই তা বৈধ (\lat{valid})।

\textbf{৫. বিনয়ী মিশন (\lat{mission}), দৃঢ় মুক্তি।}
মুমিন সাধারণ জীবনে বিনয়ী (\lat{humble}), কিন্তু সত্যের পক্ষে দৃঢ়। এই ভারসাম্য (\lat{balance}) ইসলামের চরিত্র (\lat{character})।

\medskip\hrule\medskip

\subsection*{৮. জীবনে প্রয়োগের পদ্ধতি}

\begin{enumerate}
\item \textbf{সাধারণ জীবনে বিনয় (\lat{humility}):} হাঁটা, পোশাক, কথাবার্তায় অহংকারের ভঙ্গি এড়িয়ে চলুন — কারণ কিবর জান্নাতের পথ বন্ধ করে।
\item \textbf{সত্যের পক্ষে দৃঢ়তা (\lat{firmness}):} অন্যায়ের সামনে নত না হয়ে, সত্য বলতে দৃঢ় থাকুন — এটি যুদ্ধের দৃপ্ততার মতোই আল্লাহর পছন্দ (\lat{favourite})।
\item \textbf{নিয়ত পরীক্ষা করুন (\lat{self-reflection}):} কোনো কাজের আগে জিজ্ঞেস করুন — এটা কি আল্লাহর জন্য, নাকি মানুষের কাছে বড় দেখানোর জন্য?
\item \textbf{নবীজির নজরদারি অনুশীলন (\lat{practice}):} নিজের চলাফেরা-আচরণকে নিয়মিত যাচাই (\lat{check}) করুন — যেন অহংকারের ইঙ্গিতও না থাকে।
\item \textbf{বীরত্ব ও অহংকারের পার্থক্য (\lat{difference}):} নিজের জ্ঞান, শক্তি বা সম্পদ দিয়ে মানুষকে হেয় করবেন না; কিন্তু সত্যের পক্ষে কখনো দুর্বল (\lat{weak}) হবেন না।
\end{enumerate}
\medskip\hrule\medskip

\subsection*{৯. রেফারেন্স}

\begin{itemize}
\item আল-মু'জামুল কাবীর (তাবরানী), হাদিস ৬৫০৮
\item সুনানে ইবনে মাজাহ, হাদিস ৪১৬৪ (হাসান)
\item সহীহ মুসলিম, হাদিস ৯১; সহীহ বুখারি, হাদিস ৫৭৮৮
\item সূরা আল-আনফাল, আয়াত ৬৫; সূরা মুহাম্মদ, আয়াত ৪
\item উহুদ যুদ্ধের প্রেক্ষাপট: তলোয়ার গ্রহণের ঘটনা (মুসলিম ২৪৭০)
\end{itemize}

\end{document}
